% TOP_PROBLEM: {"id": 3198, "title": "Acyclic edge-colouring", "category_id": 3, "category": {"id": 3, "name": "graph_theory", "display_name": "Graph Theory", "description": "Problems involving graphs, networks, and their properties.", "slug": "graph-theory", "order_index": 3, "created_at": "2026-07-31T15:26:25.671Z"}, "status": "open", "problem_number": "OPG-145", "set_id": 12}
% FIRST_POSTED: 2026-10-01
% ATTEMPT_STATUS: unresolved
% UnsolvedMath ID 3198; source catalogue code OPG-145
% !TeX program = lualatex
% GPT-6 Astra Ultra research attempt, 2026-10-01; independently checked partial work.
% Completion estimate: no numerical percentage assigned; see final status and literature audit.
\documentclass[11pt,a4paper]{article}
\usepackage[margin=25mm]{geometry}
\usepackage{fontspec}
\setmainfont{lmroman10-regular.otf}[BoldFont=lmroman10-bold.otf,
 ItalicFont=lmroman10-italic.otf,BoldItalicFont=lmroman10-bolditalic.otf]
\usepackage{amsmath,amssymb,mathtools}
\usepackage{unicode-math}
\setmathfont{latinmodern-math.otf}
\AtBeginDocument{\let\setminus\smallsetminus}
\usepackage{xurl,hyperref}
\hypersetup{colorlinks=true,urlcolor=blue}
\setcounter{secnumdepth}{0}
\setlength{\parindent}{0pt}
\setlength{\parskip}{5pt}
\setlength{\emergencystretch}{3em}
\begin{document}
\section*{Acyclic edge-colouring}\label{3198}
{\small UnsolvedMath ID 3198 · OPG-145 · Graph Theory · OpenGarden}
\subsection{Definitions and mathematical statement}
Let $G=(V,E)$ be a finite simple undirected graph,
and put $\Delta=\max_{v\in V}\deg_G(v)$, with
$\Delta=0$ for the empty graph. A proper
$q$-edge-colouring is a function
$c:E\to\{1,\ldots,q\}$ assigning different
colours to any two edges sharing an endpoint.
Such a colouring is acyclic if every simple
cycle in $G$ uses at least three colours.

Equivalently, for every distinct pair of
colours $a,b$, the spanning subgraph with
edge set $c^{-1}(\{a,b\})$ is a forest.
Properness forces a two-coloured cycle,
if present, to alternate its colours and
therefore have even length. Eliminating
all such cycles is the additional acyclic
requirement.

The conjecture states that every finite
simple graph has an acyclic proper edge
colouring using at most $\Delta+2$ colours:
\[
 \exists c:E\to\{1,\ldots,\Delta+2\}
 \quad c\text{ proper and every cycle uses at least three colours}.
\]
Unused colours are allowed. The bound uses
the maximum degree of the original graph,
and is independent of its order, girth,
or connectedness.

The word acyclic refers to the union of
each pair of colour classes, not to the
whole graph, which may contain many cycles.
A proper edge colouring already makes
each individual colour class a matching;
the difficulty is controlling their pairwise
unions simultaneously with only two colours
beyond the degree bound.
\subsection{Short English statement}
Can every simple graph be properly edge-coloured with at most maximum degree plus two colours so that no cycle uses only two colours?
\subsection{Research attempt}
\paragraph{Scope and process.}
GPT-6 Astra Ultra, 1 October 2026. The catalogue formulation is retained above. This attempt checks the original problem and primary literature, proves the restricted claims stated below, and identifies the exact limit of its conclusions. Reproved facts are not claimed as novel results.

\paragraph{Literature and attack.}
Esperet and Parreau [S3] prove a general $4\Delta-4$ bound using entropy compression. This checked result supplies context, not the desired $\Delta+2$ bound, and is not asserted to be the current best bound. We prove the catalogue bound for cactus graphs by a direct extension construction and establish an exact sharp example $K_4$. An acyclic colouring must be proper on edges and must give every cycle at least three colours.

\paragraph{Extension across the blocks of a cactus.}
A cactus is a graph whose blocks are single edges or simple cycles; equivalently, any two cycles have at most one vertex in common. In each connected component root the block tree and process its blocks outward. Each newly processed block meets the coloured part in at most one attachment vertex $v$; all its other vertices are new. Isolated vertices require no colour. Fix a palette of size $q=\Delta+2$.

For a single-edge block, choose a colour unused by the already coloured edges at its attachment vertex. There are at most $\Delta-1$ such edges, so a colour is available. Suppose instead that the new block is a cycle of length $\ell\ge3$. Two of its edges meet $v$, and at most $\Delta-2$ previously coloured edges do. Choose two distinct unused colours $a,b$ for those two cycle edges. Since a graph containing a cycle has $\Delta\ge2$, its palette has at least four colours.

Number the cycle edges in cyclic order $e_1,\ldots,e_\ell$, with $e_1,e_\ell$ incident to $v$, and fix their colours as $a,b$. If $\ell=3$, colour $e_2$ with any colour outside $\{a,b\}$. If $\ell\ge4$, give $e_2$ such a third colour $c$. Colour $e_3,\ldots,e_{\ell-2}$ successively, each different from its predecessor. Finally colour $e_{\ell-1}$ differently from both its predecessor and $b$; at most two forbidden colours leave a choice. For $\ell=4$ the middle range is empty and the final step colours $e_3$, as intended.

This gives a proper edge-colouring of the new cycle, uses at least three colours on it, and creates no conflict at the attachment vertex. Later blocks use colours unused at their own attachment vertices, preserving properness. Every simple cycle of a cactus is one of its cycle blocks: a simple cycle cannot pass through a cutvertex into another block and return without repeating that vertex. Hence no additional cycle can become bichromatic across several processed blocks. The construction proves the $\Delta+2$ bound for every finite cactus, including disconnected ones.

\paragraph{Why two spare colours can really be needed.}
For $K_4$, $\Delta=3$ and every colour class in a proper edge-colouring is a matching of size at most two. Six edges coloured with at most four colours require at least two classes of size two; otherwise there would be at most $2+1+1+1=5$ edges. Each two-edge matching in $K_4$ is perfect. Two disjoint such classes therefore form a four-cycle using only their two colours. Thus no proper four-edge-colouring of $K_4$ is acyclic.

Five colours suffice: give the disjoint edges $12,34$ one colour, and give each of $13,14,23,24$ a different new colour. This is proper. A two-coloured cycle in a proper colouring alternates its two colours and has at least two edges of each; here only one colour occurs twice. No such cycle exists. We have proved
\[
 a'(K_4)=5=\Delta(K_4)+2.
\]
The exact equality checks the additive constant in the catalogue rather than merely producing a loose upper bound.

\paragraph{Verification and first missing step.}
The script \texttt{scripts/check\_attempt\_batch02\_graphs\_2026\_10\_01.py} checks the explicit five-colouring and exhausts all assignments of four colours to the six edges of $K_4$, rejecting improper assignments and checking the two-colour subgraphs of every proper one. It also checks the cycle-extension rule for lengths three through twelve and all distinct endpoint-colour choices with palettes of size four through eight.

For an arbitrary graph, a new edge or path can close several cycles involving already coloured edges. Avoiding a two-coloured cycle in one location can conflict with avoiding another, and the block-tree induction no longer isolates the constraints. The cactus proof does not supply a recolouring invariant for overlapping cycles. The $K_4$ obstruction further shows that simply asking for a proper $\Delta+1$ colouring cannot meet the target's acyclicity requirement.

\paragraph{Final status.}
The full bound is proved for cactus graphs, and $K_4$ proves sharpness of its general proposed constant. The assertion for every finite simple graph is \textbf{UNRESOLVED} by this attempt.

\subsection{Sources}
\begin{itemize}
\item[S1] ULAM.AI and the UnsolvedMath Contributors, supplied problems.json, record 3198 (OPG-145), OpenGarden collection. Catalogue identity and source transcription; CC BY 4.0.
\url{https://huggingface.co/datasets/ulamai/UnsolvedMath}.
\item[S2] Open Problem Garden, Acyclic edge-colouring. Original problem and discussion; the mathematical formulation below is expanded from the supplied source record.
\url{http://www.openproblemgarden.org/op/acyclic_edge_coloring}.
\item[S3] Louis Esperet and Aline Parreau, \emph{Acyclic edge-coloring using entropy compression}, arXiv:1206.1535v3 (22 February 2013), European Journal of Combinatorics 34 (2013), 1019--1027. Checked general-bound context.
\url{https://arxiv.org/abs/1206.1535v3}.
\end{itemize}
\end{document}
